\vspace{2em}

在守旧联盟统治区域的东北部，有一座重要的中等城池“抚远”。

此城连接数条重要商路，是向东北方向诸多前哨、矿点转运物资的关键枢纽，城高池深，阵法齐全。常驻于此的超过五百名守军，由紫阳宗、丹鼎阁及其附属宗门的修士混编而成，统兵城守更是紫阳宗一位资历不浅、修为达到玄脉境融贯期的修士，赵镇岳。

在遥控四方的陈炎看来，叶牧谦及其主力被死死钉在云霞山，抚远城地处后方，守备充足，将领可靠，乃是绝无可能出问题的稳固支点，是供应前线血脉上不可或缺的一环。

他忽略了一个朴素至理：再高的城墙，再深的护城河，终究是由一块块砖石垒砌，由一个个活生生的人驻守。而砖石会风化，人心，更会变。

赵镇岳年近五旬，面容方正，眉宇间带着经年军旅生涯留下的风霜与刻痕。他的修为在玄脉境融贯期停滞多年，大道前路已渺，这才被宗门派来担任这安稳、职责分明但远离权力核心的城守之职。他性格里有旧式军人的一丝耿直乃至迁腐，对麾下兵卒和城中百姓谈不上多么体恤仁慈，但至少遵循着一套在他看来天经地义的规矩：该发的饷银尽力筹措，该收的税赋照章办理，不似某些同门那般，将手中权柄视为肆意盘剥、中饱私囊的工具。

正因这份尚未完全泯灭的底线，他对近年来宗门——尤其是陈炎一系——为推动战争而越发酷烈无情的举措，内心积郁着越来越多难以言说的微词。他亲眼见过附属宗门的弟子如何借搜查“问道途奸细”之名，行敲诈勒索散修之实；见过丹鼎阁的商队如何凭借垄断地位，肆意压低本地药农的收购价，逼得人家破人亡。更让他夜深人静时感到隐隐不安的，是麾下那些出身寒微的士卒眼中，日益浓重、几乎化不开的麻木，以及麻木之下偶尔闪过的、对不公命运的愤懑之火。

转变的契机，始于半年前一次例行公事般的剿匪。

有探报称，一伙问道途乱党流窜至抚远城西面的莽莽山林，劫掠了一支运送“慰劳物资”前往某处前哨的小型车队。赵镇岳奉命率两百士卒出城清剿。

战斗过程近乎一场闹剧，那伙“乱党”人数寥寥，战力不明，甫一接触，便远遁山林，只留下些无关紧要的杂物和早已冰凉的篝火痕迹。

就在赵镇岳心中暗骂情报不准、准备收兵回城时，麾下一名负责侧翼侦查的什长——一个名叫韩铁、同来自紫阳宗、平素沉默寡言但办事稳妥的青年——带着两名手下，押着一个受伤被擒的年轻人来到他面前。

“大人，抓到一个落单的，”韩铁的声音压得很低，带着一贯的沉稳，“看其服饰举止……不像寻常山匪流寇。”

赵镇岳目光扫向那被俘者。

对方很年轻，约莫二十出头，衣衫是寻常散修穿的粗布材质，却浆洗得干净整洁，身上有几处新添的皮肉伤，鲜血浸染了布料，但其神情却不见多少惊慌失措，尤其那双眼睛，清澈亮澈，在逆境中仍保持着一种奇异的沉静。

最让赵镇岳心下惊疑的是，以他玄脉境的感知，竟在此人身上察觉不到丝毫灵枢径修士特有的灵力波动，反而有一种绵长沉凝、迥异于此世主流却又生机勃勃的气息，在其体内隐隐流转——这与宗门卷宗里对问道途修士那模糊却特征鲜明的描述，隐隐吻合。

赵镇岳挥手屏退左右亲兵，只留下似乎完全可信的韩铁在场，沉声喝问：“你是问道途的人？”

年轻人抬手擦去嘴角渗出的血丝，动作不见狼狈，坦然迎上赵镇岳审视的目光：“是。在下林河，云隐别院记名弟子。”

“好大的胆子！”赵镇岳须发微张，玄脉境修士的威压自然流露，“区区记名弟子，也敢劫掠我紫阳宗的车队？！”

林河闻言，却轻轻笑了起来，那笑容里没有猖狂，反而带着一种让赵镇岳颇为刺目的了然与讽刺：“车队？

“大人说的是那几辆满载着‘醉仙酿’、‘灵犀脯’，专供前方大营某几位长老享用的车队么？”

他顿了顿，目光变得更加锐利，“赵城守，前方将士在云霞山下日日浴血，听闻连最基础的疗伤丹药都时常短缺，这些所谓的‘慰劳品’，究竟是去慰劳谁的？”

赵镇岳喉咙一哽，竟一时语塞。

车队里装的是什么，他岂会不知？

宗门上层某些人的奢靡做派与前线艰困形成的刺眼对比，早已不是秘密。

但他身处其位，有些事看破不能说破。

他只能把脸一板，厉声道：“巧舌如簧！你们这些乱党，四处煽风点火，破坏后方安定，才是致使生灵涂炭的罪魁祸首！”

林河缓缓摇头，语气依旧平静，却字字如有千钧：“我们从未煽动任何人。我们只是将血淋淋的事实，将另一种活法的可能，摆在了那些被压迫、被蒙蔽的人面前。

“赵城守，你坐镇抚远多年，耳目通达。你可曾听闻，我问道途义军所到之处，主动劫掠过一户平民，欺辱过一名无辜散修？可曾见过我们攻占某地后，如贵方某些军爷那般，纵兵抢掠三日，或是强征暴敛、竭泽而渔？”

赵镇岳再次陷入沉默。

他确实听到过一些截然不同的风声。从那些商人、游侠零散的交谈中，从一些难以完全封锁的渠道里，他隐约得知，问道途义军行事极有章法，只针对联盟官仓和死硬分子，对平民百姓甚至堪称秋毫无犯，有时还会将缴获的粮食物资分发给穷苦人。这与宗门宣传中那烧杀抢掠、无恶不作的邪魔外道形象，相差何止万里。

见赵镇岳脸色变幻，林河知道话语已触及对方心防，便说出了那直指本心的一句：“《问道戒律》开宗明义，‘仰不愧于天，俯不怍于人’。我等所求，无非是开辟一条道路，让无数无根无基的散修，让终年劳作却不得温饱的百姓，也能凭借自身的汗水与毅力，存续性命，活得有几分人的尊严。”

林河的目光深深看进赵镇岳眼中，“请您扪心自问，在这世道，在这抚远城内城外，还有多少人，能堂堂正正地说自己活得‘不愧不怍’？您自己呢？”

“仰不愧于天，俯不怍于人”！

这十个字，如同一声暮鼓晨钟，又似一根烧红的钢针，狠狠刺入赵镇岳灵魂深处某个早已被世俗尘埃覆盖、却从未真正死去的角落。

他想起了年少时初入山门，也曾怀揣过斩妖除魔、护卫一方的热血梦想；想起了这些年来目睹的越来越多无力改变的龌龊与不公；更想起了麾下那些朴实士卒，在收到家书得知赋税又涨时，眼中那无法掩饰的忧虑与对前途的茫然。

复杂的情绪在胸中翻涌，最终，赵镇岳没有回答林河的问题，只是有些疲惫地挥了挥手，对侍立一旁的韩铁道：“给他包扎伤口，押回城中……单独关押于僻静处，不准用刑。等我亲自发落。”

回到抚远城，赵镇岳的心绪再也无法平静。

几日后的一个夜晚，他屏退左右，再一次独自去见了林河。

这一次，林河没有慷慨激昂地陈述理念，反而像一位关心民生疾苦的访客，细致地问起了抚远城的现状：今年粮价几何？各类税赋名目有多少，底层可能承受？守城士卒的饷银能否按时足额发放？城中散修坊市的生意是否兴旺，可有受到盘剥？

这些问题，有的赵镇岳清楚，有的却只能含糊以对——他虽为城守，但对底层具体生态的了解，远不如一个真正深入其中的观察者。

林河便不急不缓，向他讲述起其他类似规模城池，在守旧联盟治下与在问道途影响下的真实对比。数据详实，事例具体，何人因何破产，何处因何得以喘息，言之凿凿，逻辑清晰，由不得赵镇岳不信。

他第一次如此直观而深刻地意识到，问道途和他们一样，同样有着清晰的理念、严明的纪律，甚至是对未来如何治理一方、改善民生的具体构想。相比之下，宗门上层某些人穷兵黩武、竭泽而渔的做法，显得如此短视、粗暴，且不得人心。

与此同时，赵镇岳察觉到，自己颇为信任的部下韩铁，在汇报看守情况或谈及城中琐事时，也常常会“无意”地透露一些信息：某个老兵的家眷因交不起新加的“防务捐”被税吏打了；某个相识的散修朋友最近精神头好了不少，私下透露在练一套“奇特的养气法门”；城中酒肆茶楼里，对宗门久攻云霞山不下、反而赋税越来越重的抱怨声，正在悄然滋长，虽不敢明目张胆，却如地底暗流，涌动不休。

赵镇岳不是愚钝之人。他渐渐明白，韩铁恐怕也早已被渗透，甚至可能就是对方早早埋下的一枚棋子。

但他没有点破。

一方面，韩铁行事依旧稳妥可靠，并未做出任何损害他本人或城池防务的举动；另一方面，那些通过韩铁之口传来的消息，虽然刺耳，却往往是真实存在的、被他有意无意忽略的城中疾苦。一种微妙而奇异的默契，在赵镇岳、林河、韩铁这三个身份迥异的人之间，悄然形成。

导火索在三个月后被点燃。陈炎为应对日益猖獗、防不胜防的“后方袭扰”，决心强化清剿，命令各后方城池加征一批紧急物资，并抽调部分守军补充前线损耗。命令传到抚远，要求赵镇岳抽调一百名修士即刻前往报到，同时上缴城库中灵石储备的三分之一、疗伤丹药的一半。

这道命令，成了压垮赵镇岳心中那架摇摆天平的最后一块巨石。

抽调百名修士，几乎等于抽走了抚远城三分之一的机动防御力量，城池安全顿时空虚；而上缴大量关键物资，不仅动摇了守城的物质基础，更意味着本已紧张的城内储备将急剧缩水，随之而来的必是物价飞涨，民不聊生，怨声载道。

赵镇岳在城守府的书房内独坐了一整夜。案头那纸调令冰冷而沉重，窗外是抚远城沉睡的轮廓。林河关于民心向背的论述，韩铁平日透露的点点滴滴，宗门上层的苛刻短视，与这座他驻守多年、虽无深情却已有责任的城池可能面临的混乱破败未来，交织在一起，反复灼烧着他的理智。

天光将亮未亮之际，书房内的烛火燃尽了最后一滴油脂，悄然熄灭。在黎明前最深的黑暗中，赵镇岳抬起了头，眼中最后的犹豫已然散去，取而代之的是一种孤注一掷的决绝。

他秘密召来了韩铁。没有迂回，没有试探，声音因一夜未眠而沙哑，眼神却亮得惊人：“联系你们的人。”

韩铁身躯微不可察地一震，瞬间明白了这句话的含义，眼中爆发出压抑已久的、惊人的神采。

但他依然强行克制住激动，确认般低声问：“您指的是……紫阳宗方面？”

“当然不是紫阳宗。”

赵镇岳斩钉截铁，随即清晰地道出了自己的条件，“抚远城……可以易帜。但我有三个条件：第一，不得滥杀任何投降的守军及城中无辜百姓；第二，城中所获库藏，必须优先用于安抚民生、救治伤患；第三，我赵镇岳个人之生死荣辱，无足轻重，但请务必放过我家中老小，保他们平安离去。”

韩铁再无犹豫，单膝跪地，抱拳行礼，声音因激动而微颤：“赵城守深明大义！您所提条件，正是问道途一贯秉持之道义所为！韩铁以性命担保，必即刻将城守大义与条件，如实禀报白先生、沈姑娘！”

联络通过韩铁掌握的、深埋在敌后网络中的隐秘渠道，以最快速度传递出去。五日后的深夜里，一只毫不起眼、却附有特殊灵力印记的纸鹤，穿过层层阵法防护，悄然落入赵镇岳手中。他展开纸鹤，上面只有一行蝇头小楷，墨迹沉稳：

“七日后，子时三刻，东门。白、沈。”

赵镇岳将这短短一行字反复看了数遍，仿佛要将每一个笔画都刻入脑海。然后，他指尖腾起一缕火焰，将纸鹤点燃，看着它化为灰烬，飘散在夜风中。他知道，退路已断，从这一刻起，他要么与这座城一起获得新生，要么便坠入万劫不复的深渊。

接下来至关重要的七天，赵镇岳仿佛在刀尖上行走。

他必须表现得一切如常，甚至更加勤勉。

面对宗门派来催促物资和兵员的使者，他做出了恰到好处的为难与抱怨，以“清点库藏需要时间”、“士卒调动涉及防务交接”等理由，成功地将时间拖延了下来。与此同时，他暗中以“整备防务”、“预防奸细”为名，进行了一系列至关重要的布置：将库房中部分关键物资转移到更隐秘或更易控制的地点；将几个要害岗位的守卫，尤其是东门附近的哨位，逐步换上了经过韩铁筛选、绝对可靠或已被争取过来的士卒；对少数几个紫阳宗、丹鼎阁的死忠分子及其掌控的小队，则拟定了详细的监控与制伏计划。

整个抚远城，表面上依旧秩序井然，运转如常，但在平静的水面之下，一股决定城池命运的暗流，已经蓄势待发。

第七日，夜幕早早降临，厚重如墨，无星无月，正是改天换地之时。

子时将近，抚远城仿佛沉入酣睡，只有城墙上的气死风灯在夜风中孤零零地摇曳，昏黄的光晕照亮了值守士卒们略带困倦的脸庞。东门城楼上，今夜的值守官，正是什长韩铁。他像往常一样，带着十余名亲信士卒，神情严肃地巡视着城防。

时间一点一滴流逝，指向那个约定的时刻。

子时三刻整。

韩铁脚步沉稳地走到城楼一角，那里看似随意地堆放了一些用于夜间照明和烽火传讯的干燥柴薪、油脂等引火之物。他警惕地扫视四周，确认无人异常关注此处，随即迅速从怀中掏出火折，迎风一晃，幽蓝的火苗窜起，被他果断地引向柴堆。

“嗤啦——！”

干燥的物料遇火即燃，火舌猛地蹿起，越烧越旺，橘红色的烈焰在漆黑如铁的夜幕映衬下，显得格外刺目而耀眼！

“走水了！东门楼走水了！”几乎是立刻，城头响起了“惊慌失措”的喊叫，警锣被奋力敲响，刺耳的“铛铛”声划破夜的寂静，迅速传遍附近城墙。城头守军一阵骚动，附近几座哨塔的视线和注意力，全被这突如其来的“意外”火灾吸引了过去。

然而，这烽火，并非意外，而是号令！

几乎就在东门楼火光亮起的同一刹那，在抚远城外数里处一片早已埋伏多日的密林中，数百双在黑暗中依旧清亮的眼睛，同时睁开！

以白言冰、沈瑶为首，吴刚、陈默等核心骨干协同，超过三百名历经血火淬炼的问道途精锐游击队战士，以及闻讯响应、自愿前来参战的附近区域散修、乃至部分被成功策反的小家族私兵，总数近五百人，如同蛰伏已久的猛兽，骤然跃出藏身之地！

急促而整齐的脚步声、兵刃与轻甲摩擦的沙沙声，汇成一股低沉却充满压迫感的洪流，向着洞开的城门方向（他们确信那里将会洞开），发起了无声而迅猛的冲锋！

队伍最前方，白言冰一袭朴素的青色劲装，手握长剑，身影如电，眼神沉静如古井寒潭，不见波澜，却蕴含着劈开一切阻碍的决心。

紧随其侧的沈瑶，黑衣如墨，身形在夜色中几乎融入背景，唯有那双眼眸，在疾驰中闪烁着千幻流光般冰冷而锐利的光芒，仿佛已锁定了城中所有需要拔除的顽抗节点。他们，以及他们身后每一个战士，为这一天，已经等待了太久，积蓄了太久。

城头，真正的混乱与更替，在火光与喊叫的完美掩饰下，高效而致命地展开了！

在火灾吸引注意力的瞬间，韩铁和他那十余名早已准备就绪的心腹，如同捕猎的豹群般突然发动！他们目标明确，配合默契，以迅雷不及掩耳之势，扑向东门附近少数尚未反应过来的、已知的死忠紫阳宗、丹鼎阁修士及其铁杆附庸。

刀光剑影在火光映照下短暂闪烁，闷哼与短促的惨叫声被更大的嘈杂掩盖，反抗在组织性的突袭面前迅速瓦解。

“控制门闸！开城门！”韩铁低吼一声，挥刀砍断粗重门闩上的铁锁，与几名膂力惊人的壮汉合力，推动那扇厚重的包铁城门。

“吱嘎嘎——轰隆！！”

城门洞开的沉重声响，如同惊雷，在寂静的夜晚传出去极远，也传入了城外汹涌而来的洪流耳中。

“城门已开！随我进城！按预定计划，夺取要点，镇压顽抗，保护百姓！”白言冰的命令穿透夜色，清晰地传入每一个冲锋战士的耳中。

随即，白言冰一马当先，身形化作一道青色流光，瞬息间掠过护城河，穿入洞开的城门。沈瑶如影随形，刀光在她手中仿佛拥有了生命，每一次闪烁，都精准地带走一名试图组织抵抗的士兵的性命，为后续大军扫清道路。

数百精锐如决堤之水，汹涌灌入抚远城，随即按照早已演练过无数次的方案，迅速分成数十支精干小队，扑向各处战略要害：城守府、中心武库、传送阵法枢纽、各大官仓、监狱以及通往其他城门的交通要道。

与此同时，城内的内应行动也全面爆发。赵镇岳亲自坐镇，以“救火”、“有大量匪徒趁乱潜入城内”为名，调动了自己绝对掌控的亲兵队，以及被韩铁等人长期渗透、拉拢的绝大部分中下层军官和士卒，向那些驻扎在军营、武库等地的死硬分子发动了内外夹击的突袭。

许多紫阳宗、丹鼎阁的修士在睡梦中被同僚的刀剑指住咽喉，或是在执勤时被身后“战友”突然反水，大多茫然失措，零星抵抗在早有准备的镇压下迅速熄灭。

夺取抚远城的战斗，激烈却短暂。反抗主要集中于几处据点，在问道途精锐与起义守军的联手打击下，未能形成有效抵抗。天色将明未明之际，城中主要区域已基本被控制。

白言冰与沈瑶亲自带队，直扑城守府。在那里，他们见到了早已卸去甲胄、只着一身整洁常服，独自立于庭院中央等待的赵镇岳。他脚下，横放着自己那柄伴随多年的佩剑。

见白言冰、沈瑶到来，赵镇岳抱拳，深深一礼，声音带着一夜未眠的疲惫，却异常清晰：“罪将赵镇岳，约束部下不力，致使城池易手，愿领一切责罚。只求……只求贵军信守前约。”

白言冰快步上前，郑重拱手还礼，语气诚恳而有力：“赵城守深明大义，于紧要关头弃暗投明，不仅使无数将士百姓免遭无谓死伤，更保全此城元气，功莫大焉！我白言冰，谨代表云隐别院叶牧谦先生，欢迎赵城守加入我问道途行列，共襄义举！此前约定诸事，我问道途必一字不易，全力履行！”

沈瑶亦微微颔首，清冷的容颜上罕见地露出一丝温和：“赵城守能以满城生灵为念，沈瑶感佩。眼下城中初定，百废待兴，诸多秩序维系、人心安抚事宜，还需赵城守鼎力相助。”

赵镇岳听闻此言，心中那块高悬已久、重于千钧的巨石，轰然落地。一股混合着如释重负、愧疚与新生希望的热流冲上眼眶，他强行忍住，再次重重抱拳，声音微哽：“镇岳……敢不从命！定竭尽所能，助两位稳定抚远，戴罪立功！”

而韩铁与白言冰两人相遇时，更是相互抱拳一礼。

两年半之前，青鸾秘境之中，救命之恩，谁敢忘记？谁又能忘记？！

当第一缕晨光刺破东方的云层，照耀在抚远城头时，一面崭新的旗帜，在万众瞩目下缓缓升起，取代了那面象征紫阳宗权威的烈焰战旗。

旗帜底色是洁净的素白，中央以遒劲的笔墨书写着十个大字——“仰不愧于天，俯不怍于人”。在这核心信念之下，又分三行列着更为具体、直指行止的律条：

“一切行动听指挥”

“不拿群众一针一线”

“一切缴获要归公”

这，便是叶牧谦根据前世宝贵经验，结合此世实际，为这支新生人民军队确立的“三大纪律”。它简单，直白，却蕴含着超越所有宗门戒律的、与最广大民众血肉相连的深刻力量。

城中百姓在胆战心惊中推开家门，预料中的烧杀抢掠并未发生。街道上虽然有列队巡逻或匆匆往来的士卒，其中许多还是他们熟悉的本地守军，但秩序井然。

更令人惊异的是，街头巷尾的显眼处，一夜之间贴满了墨迹未干的安民告示。告示以浅显易懂的文字宣布：即刻废除守旧联盟在抚远城推行的一切苛捐杂税；打开官仓，以极平价出售粮食物资，对鳏寡孤独及赤贫者进行赈济；城中所有伤患，无论军民，皆可前往指定地点，获得问道途医修的免费救治。

而最让城中数量庞大的散修群体心跳加速、热血上涌的，是告示最后一条：即日起，于城中原演武场设立“讲道坪”，每日由问道途修士公开讲授《养气篇》基础法诀与修行疑难，不问出身，不论来历，皆可前来听讲！

疑虑、恐惧、观望……在随后几个时辰里，迅速被一件件实实在在发生的事情碾碎。当第一个大胆的贫苦老人真的用几枚铜钱从官仓买到足以果腹的灵米，当第一个在昨夜零星冲突中受伤的孩童被随军的医修温柔包扎、喂下药散，当第一个散修怀着朝圣般的心情踏入“讲道坪”，听到那深入浅出、直指大道的讲解时……积蓄已久的民心，如同冰封的河面在春日暖阳下轰然迸裂，化作汹涌澎湃的春潮！

“问道途万岁！”

“叶先生万岁！”

“白将军！沈仙子！”

欢呼声起初在几个街区零星响起，随即如同野火燎原，迅速蔓延至全城。那不是被胁迫的呼喊，而是发自肺腑的、对新生与希望的狂热礼赞！

白言冰、沈瑶等人深知，夺取城池只是第一步，巩固胜利果实更为关键。他们迅速清点战果：缴获的灵石、各类矿石、药材、粮食堆积如山，数量之巨，不仅足以支撑云霞山前线数月之用，更能极大缓解己方各根据地的物资压力。另一面，吴刚的后勤网络很快跟进，抚远城从一个孤立的堡垒，瞬间转变为连接云霞山与广袤敌后农村根据地的战略支点，物资人员流转效率陡增，真正实现了农村与城市的联动，进可攻，退可守。

然而，抚远城易帜的最大意义，远不止于军事和物资层面。它成了一个极具冲击力和示范效应的宣传样板。白言冰命人将战役概况、入城后的安民措施、新旧治理的鲜明对比，通过早已枝繁叶茂的地下网络，多渠道、全方位地迅速散播出去。

消息所至，石破天惊！

那些仍在守旧联盟统治下苦苦支撑、观望犹豫的中小城池守军，那些饱受压迫却不知出路在何方的散修民众，第一次如此清晰地看到：原来反抗并非必死，易帜竟能带来更好的生活；原来那面“仰不愧天，俯不怍人”的旗帜，不仅是一句口号，更能实实在在地庇护一方，带来公平与希望。守旧联盟后方本就因“农村包围城市”而暗流汹涌的人心，被“抚远事件”这把烈炬彻底点燃，彻底沸腾！

更多的“韩铁”在黑暗中更加活跃地串联，更多的“赵镇岳”在内心的挣扎与权衡中，悄然倾向另一边。

抚远一夜，乾坤始定。

这座中型枢纽城池的易手，如同在守旧联盟看似铁板一块的统治版图上，用最猛烈的方式凿开了一个巨大的、鲜血淋漓的缺口。光芒与希望从这个缺口汹涌灌入，而恐惧与动摇则从此处向内疯狂蔓延。

这不再仅仅是一次战术层面的奇袭胜利，而是一次政治与战略上的决定性转折，标志着问道途从长期的战略防御与农村渗透，正式迈入了武装夺取政权、巩固扩大根据地的战略反攻阶段。

内域道争的天平，经此一役，倾斜的速度骤然加剧，再难逆转。

消息通过特殊渠道，以最快速度传回巍峨屹立的云霞山。当叶牧谦展开密信，看到“抚远已定，白、沈、赵等皆安，人心归附，三大纪律张布于城头”等字样时，他久久伫立，遥望东北方向，脸上露出了自战争爆发以来，最为舒展、如释重负的欣然微笑。

那笑容里，有对弟子们成长的欣慰，有对策略成功的印证，更有对那不可阻挡的历史车轮滚滚向前的深切感悟与感慨。

而与此同时，紫阳宗内，接到十万火急、损失详报的陈炎，在短暂的死寂之后，猛然暴起，将他平日最珍爱的一套千年暖玉茶具狠狠掼在地上，摔得粉碎！瓷片与茶水四溅，映照着他那张因极度震怒、挫败以及一丝连他自己都不愿承认的恐慌而彻底扭曲铁青的脸。

“叶牧谦……韩铁……赵镇岳……好，好得很！”他咬牙切齿，声音从喉咙深处挤压出来，带着血腥气，“你们以为，夺下一城，就能翻天覆地吗？！”

他猛地转向噤若寒蝉的传令修士，眼中闪烁着近乎疯狂的凶光，嘶声吼道：

“传我急令！前线围山部队，再给我抽调……不！急令周文彬，让他立刻放下手头一切事务，滚来见我！立刻！马上！”

风暴，将以抚远为中心，向着更广阔、更惨烈的领域，悍然席卷。真正的战略决战，已被这成功策反、武装夺取的第一座城池，推到了历史的门前。